\section{Boundaries}

The implemented claim is a comparison-audit replay claim. RCOUNT verifies that
the declared Kaplan-Markov taint product or MACRO product follows from the audit
record fields and sample-step categories. It does not prove that a jurisdiction
retrieved the correct physical ballot, that a CVR-to-ballot link was trustworthy,
or that a batch was handled under a valid chain of custody.

The method is also not ballot polling. BRAVO and Minerva consume ballot
observations as votes for a reported winner, loser, or other choice. The
Kaplan-Markov comparison surface consumes overstatement evidence derived from a
CVR or reported total and a hand interpretation of the same unit. If that link
is missing, the statistical replay can still be syntactically present but the
external audit claim is outside this contract.

The MACRO path is deliberately field-gated. RCOUNT uses the published-formula
MACRO replay only when the package declares $N$, $V$, and gamma together. Without
those fields, \texttt{kaplan-markov-comparison-v1} falls back to the older
taint-product replay surface. The project plan therefore leaves one validation
boundary open: compare the MACRO package path against external public audit
fixtures when a source exposes the ballot count, reported margin, gamma, and
ordered overstatement categories.
